% TEMPLATE-MILC-v3.tex - plantilla maestra para todas las guias MILC v3
% Ver scripts/generadores/build-guia-milc.py para la lista completa de
% claves esperadas. El script valida que no queden placeholders sin
% reemplazar antes de escribir el archivo final.

\documentclass[11pt,letterpaper]{article}
\usepackage[margin=1.75cm]{geometry}
\usepackage[table]{xcolor}
\usepackage{fontspec}
\usepackage[spanish]{babel}
\usepackage{tikz}
\usepackage[most]{tcolorbox}
\usepackage{tabularx}
\usepackage{array}
\usepackage{fancyhdr}
\usepackage{titlesec}
\usepackage{ragged2e}
\usepackage{enumitem}
\usepackage{microtype}

\setmainfont{Helvetica Neue}
\setsansfont{Helvetica Neue}
\setlength{\parindent}{0pt}
\setlength{\parskip}{6pt}
\hyphenpenalty=200
\exhyphenpenalty=200
\pretolerance=400
\tolerance=2000
\setlength{\emergencystretch}{1em}
\renewcommand{\arraystretch}{1.25}
\setlist{leftmargin=5mm,itemsep=1mm,topsep=1mm}
\newcolumntype{Y}{>{\RaggedRight\arraybackslash}X}

\definecolor{milcMagenta}{HTML}{E6007E}
\definecolor{milcVino}{HTML}{5A0038}
\definecolor{milcTurquesa}{HTML}{0093A5}
\definecolor{milcVerde}{HTML}{58A923}
\definecolor{milcMostaza}{HTML}{E5B400}
\definecolor{milcOcre}{HTML}{B86600}
\definecolor{milcNegro}{HTML}{241816}
\definecolor{milcGris}{HTML}{F7F7F4}
\definecolor{milcLinea}{HTML}{E8E2DD}
\definecolor{milcGradePrimary}{HTML}{0066FF}
\definecolor{milcGradeDark}{HTML}{003D99}
\definecolor{milcGradeSoft}{HTML}{D6E8FF}
\definecolor{milcGradeAccent}{HTML}{CFE7FF}
\definecolor{milcGradeText}{HTML}{FFFFFF}
\color{milcNegro}

\pagestyle{fancy}
\fancyhf{}
\lhead{\small Tecnología e Informática · Grado 6}
\rhead{\small Guía 22 · MILC v3}
\cfoot{\small\thepage}
\renewcommand{\headrulewidth}{0pt}

\newtcolorbox{titlebox}[1]{
  enhanced,colback=#1,colframe=#1,arc=7mm,boxrule=0pt,
  left=7mm,right=7mm,top=3mm,bottom=3mm,
  before skip=8pt,after skip=8pt,
  fontupper=\sffamily\bfseries\Large\color{white}
}

\newtcolorbox{softbox}[3]{
  enhanced,colback=#2,colframe=#1,borderline west={4pt}{0pt}{#1},
  arc=2mm,boxrule=.4pt,left=4mm,right=4mm,top=3mm,bottom=3mm,
  before skip=6pt,after skip=8pt,title={#3},coltitle=white,
  colbacktitle=#1,fonttitle=\sffamily\bfseries,fontupper=\RaggedRight,
  attach boxed title to top left={xshift=3mm,yshift=-2mm},
  boxed title style={arc=2mm, boxrule=0pt}
}

\newtcolorbox{drawbox}[2]{
  enhanced,colback=white,colframe=#1,arc=4mm,boxrule=1pt,
  left=5mm,right=5mm,top=4mm,bottom=4mm,before skip=8pt,after skip=8pt,
  title={#2},coltitle=white,colbacktitle=#1,fonttitle=\sffamily\bfseries,
  fontupper=\RaggedRight,
  attach boxed title to top left={xshift=4mm,yshift=-2mm},
  boxed title style={arc=3mm, boxrule=0pt}
}

\newtcolorbox{infoband}[1]{
  enhanced,colback=#1!8,colframe=#1!50,arc=2mm,boxrule=.5pt,
  left=4mm,right=4mm,top=2mm,bottom=2mm,before skip=4pt,after skip=4pt,
  fontupper=\small\RaggedRight
}

% Puente narrativo entre secciones (contrato editorial Conéctate).
% Da continuidad: "Acabas de X. Ahora vas a Y porque Z."
% Visualmente: banda izquierda en color del grado, texto en itálica pequeña.
\newtcolorbox{puentebox}{
  enhanced,colback=milcGradeSoft,colframe=milcGradeDark,
  borderline west={3pt}{0pt}{milcGradeDark},
  arc=0mm,boxrule=0pt,left=5mm,right=4mm,top=2.5mm,bottom=2.5mm,
  before skip=4pt,after skip=8pt,
  fontupper=\itshape\small\color{milcNegro}\RaggedRight
}
\newcommand{\puente}[1]{%
  \begin{puentebox}%
    \textcolor{milcGradeDark}{\textbf{\textrightarrow}}\hspace{2mm}#1%
  \end{puentebox}%
}

\newcommand{\linead}[1]{\noindent\color{milcLinea}\rule{#1}{0.5pt}\color{milcNegro}}
\newcommand{\respuesta}[1]{\par\vspace{2mm}\foreach \n in {1,...,#1}{\noindent\color{milcLinea}\rule{\linewidth}{0.5pt}\par\vspace{5mm}}\color{milcNegro}}
\newcommand{\checkbox}{\textcolor{milcGradeDark}{\fbox{\phantom{x}}}\hspace{5pt}}

\newcommand{\phasecard}[4]{
\begin{tcolorbox}[enhanced,colback=#3,colframe=#2,arc=3mm,boxrule=.5pt,height=3.05cm,valign=center,halign=center,left=1.5mm,right=1.5mm,top=2mm,bottom=2mm]
{\sffamily\bfseries\fontsize{22}{24}\selectfont\textcolor{#2}{#1}}\par
{\sffamily\bfseries\small #4}
\end{tcolorbox}
}

\begin{document}
\thispagestyle{empty}
\begin{tikzpicture}[remember picture,overlay]
  \fill[white] (current page.south west) rectangle (current page.north east);
  \path[fill=milcGradePrimary,rounded corners=16mm]
    ([xshift=.55cm,yshift=-.55cm]current page.north west) rectangle
    ([xshift=-.55cm,yshift=3.65cm]current page.south east);

  \node[anchor=north west,text=milcGradeText] at ([xshift=2.0cm,yshift=-1.85cm]current page.north west)
    {\sffamily\bfseries\fontsize{14}{16}\selectfont GRADO};
  \node[anchor=north west,text=milcGradeText,opacity=.82] at ([xshift=2.0cm,yshift=-2.75cm]current page.north west)
    {\sffamily\bfseries\fontsize{9}{11}\selectfont SERIE GUÍAS · TECNOLOGÍA E INFORMÁTICA};

  \begin{scope}[shift={([xshift=-3.05cm,yshift=-2.55cm]current page.north east)}]
    \fill[white,opacity=.80] (-.62,-.52) rectangle (.62,.52);
    \draw[milcGradeText,line width=1pt] (-.62,-.52) rectangle (.62,.52);
    \fill[milcGradeDark] (-.42,-.38) rectangle (-.22,.06);
    \fill[milcGradeAccent] (-.10,-.38) rectangle (.10,.24);
    \fill[milcGradePrimary!60!black] (.22,-.38) rectangle (.42,.38);
  \end{scope}

  \node[anchor=west,text=milcGradeText] at ([xshift=1.9cm,yshift=-12.85cm]current page.north west)
    {\sffamily\bfseries\fontsize{124}{124}\selectfont 6};
  \draw[milcGradeText,line width=1.7pt]
    ([xshift=4.52cm,yshift=-11.68cm]current page.north west) circle (.13cm);

  \node[anchor=west,text=milcGradeText,text width=8.7cm,align=left] at ([xshift=10.35cm,yshift=-11.6cm]current page.north west)
    {
     {\sffamily\bfseries\fontsize{19}{22}\selectfont ¿Qué es un\\procesador de\\texto? ---\\mi primera\\página digital}\\[4mm]
     {\sffamily\bfseries\fontsize{23}{26}\selectfont Guía 22-6-TIC}\\[2mm]
     {\sffamily\fontsize{13}{16}\selectfont Período 3 · Procesador de texto e Internet}\\[5mm]
     {\sffamily\bfseries\fontsize{11}{13}\selectfont Institución Educativa Sor María Juliana}\\[1mm]
     {\sffamily\fontsize{10}{12}\selectfont Autor: Dr. Álvaro Cárdenas Orozco}
    };

  \path[fill=milcGradeSoft,rounded corners=7mm]
    ([xshift=1.35cm,yshift=.85cm]current page.south west) rectangle
    ([xshift=-1.35cm,yshift=4.25cm]current page.south east);
  \draw[milcGradeDark,line width=.65pt,rounded corners=7mm]
    ([xshift=1.35cm,yshift=.85cm]current page.south west) rectangle
    ([xshift=-1.35cm,yshift=4.25cm]current page.south east);
  \node[anchor=center,text=milcNegro,text width=17.0cm,align=center] at ([yshift=2.55cm]current page.south)
    {\sffamily\fontsize{10}{12}\selectfont Metodología MILC · Apertura ancestral · Pensamiento computacional · Triángulo Dussel-Estoicismo-Floridi\\
    \textcolor{milcGradeDark}{\bfseries Producto:} Un \textbf{documento digital titulado ``Mi primer texto digital''} con \textbf{3 párrafos} (introducción, desarrollo, cierre) sobre un tema personal que tú escojas (tu pueblo, tu materia favorita, tu mascota, lo que quieras). El documento se guarda en la carpeta correcta (siguiendo lo aprendido en G6·P2·S7), tiene nombre claro siguiendo las 5 reglas, y se imprime o se muestra en pantalla al profe. En el cuaderno: \textbf{captura impresa o boceto} del documento + \textbf{pasos que seguiste} numerados.};
\end{tikzpicture}
\mbox{}
\newpage

\section*{Apertura: ¿Qué es un procesador de texto? --- mi primera página escrita en digital}

\begin{softbox}{milcGradeDark}{milcGradeSoft}{Saber ancestral}
\textbf{Doña Mercedes era la maestra rural de la vereda La Plata, en Cartago, Valle del Cauca.} Hace 40 años, no había computadores. En su escuelita de una sola aula, doña Mercedes enseñaba a los niños del campo \textbf{a escribir}: tomar el lápiz, formar las letras, juntar las palabras, ordenar las ideas en párrafos. Para ella escribir era \emph{``ordenar el pensamiento para que otros lo entiendan''}. Repetía a sus alumnos: \emph{``Mijo, antes de escribir, pensar. Antes de juntar las palabras, ordenar las ideas. Una hoja escrita con cuidado vale más que cien escritas con prisa.''}. Esa misma sabiduría sigue viviendo, pero ahora la escritura no es solo en cuaderno: también es \textbf{digital}, en programas que se llaman \emph{procesadores de texto}. Aprender a usarlos no reemplaza el cuaderno (sigue siendo fundamental); \textbf{lo complementa}. Hoy doña Mercedes te enseña a escribir digital con el mismo cuidado con el que enseñaba a escribir en papel.
\end{softbox}

\begin{softbox}{milcTurquesa}{milcGris}{Saber contemporáneo}
Un \textbf{procesador de texto} es un programa para escribir documentos en el computador. Los más usados son \textbf{Microsoft Word} (parte de Office 365, con versión gratis estudiantil), \textbf{Google Docs} (gratis, en el navegador, guarda en Drive) y \textbf{LibreOffice Writer} (gratis, instalable). Los 3 hacen lo mismo: escribir, dar formato, guardar. La diferencia es \emph{dónde guardan} (Word en tu disco; Docs en la nube; LibreOffice en tu disco) y \emph{qué tan rico es el formato} (Word es el más completo). \textbf{Lo que no cambia}: la escritura sigue siendo \emph{pensar, ordenar, escribir, revisar}. Doña Mercedes te lo dijo. Hoy aprendes el procesador como herramienta nueva, pero el oficio sigue siendo el mismo.
\end{softbox}

\begin{softbox}{milcMostaza}{milcGris}{Pregunta-puente}
\textit{Si te pidieran escribir 3 párrafos sobre tu \emph{equipo de fútbol favorito} o sobre \emph{una tradición de tu familia}, ¿\textbf{lo escribirías en cuaderno o en computador}? ¿Hay diferencia entre los dos? ¿Cuál?}
\end{softbox}

\begin{softbox}{milcVerde}{milcGris}{Saber y saber hacer}
Al terminar la clase: (1) podrás \textbf{identificar} los 3 procesadores de texto más comunes; (2) sabrás \textbf{aplicar} los 5 pasos para crear, escribir y guardar un documento; (3) podrás \textbf{distinguir} entre escribir en computador y en cuaderno (cuándo conviene cada uno); (4) habrás \textbf{creado} tu primer documento digital de 3 párrafos guardado correctamente.
\end{softbox}

\small
\begin{tabularx}{\textwidth}{>{\bfseries}p{3.0cm}Y}
\rowcolor{milcGradePrimary!12}Autor & Dr. Álvaro Cárdenas Orozco\\
DBA & Produce contenidos digitales y valida información de fuentes web (MEN, T\&I 6°)\\
Referentes & Saber ancestral de la maestra rural · Lectura crítica · Soberanía informacional\\
Producto final & Un \textbf{documento digital titulado ``Mi primer texto digital''} con \textbf{3 párrafos} (introducción, desarrollo, cierre) sobre un tema personal que tú escojas (tu pueblo, tu materia favorita, tu mascota, lo que quieras). El documento se guarda en la carpeta correcta (siguiendo lo aprendido en G6·P2·S7), tiene nombre claro siguiendo las 5 reglas, y se imprime o se muestra en pantalla al profe. En el cuaderno: \textbf{captura impresa o boceto} del documento + \textbf{pasos que seguiste} numerados.\\
\end{tabularx}
\normalsize

\puente{Hoy haces 4 cosas: comparas los 3 procesadores, aprendes los 5 pasos para escribir y guardar, escoges tu tema, y creas tu documento real.}

\begin{titlebox}{milcGradeDark}
Ruta MILC: así vamos a aprender
\end{titlebox}
\begin{tabularx}{\textwidth}{>{\centering\arraybackslash}X>{\centering\arraybackslash}X>{\centering\arraybackslash}X>{\centering\arraybackslash}X}
\phasecard{1}{milcVerde}{milcGris}{Escuta\\[-1mm]\small Observo, escucho y pregunto.} &
\phasecard{2}{milcTurquesa}{milcGris}{Sistematización\\[-1mm]\small Pensamiento computacional.} &
\phasecard{3}{milcMagenta}{milcGris}{Praxis\\[-1mm]\small Construyo y aplico.} &
\phasecard{4}{milcVino}{milcGris}{Evaluación\\[-1mm]\small Reflexión liberadora.}
\end{tabularx}


\puente{Empezamos comparando Word, Google Docs y LibreOffice en 4 puntos.}

\begin{titlebox}{milcVerde}
Fase 1 · Escuta
\end{titlebox}

\begin{softbox}{milcVerde}{milcGris}{Escena para escuchar}
\textbf{Actividad 1 · IDENTIFICA --- Comparar los 3 procesadores} (10 min · individual). En el cuaderno, dibuja una tabla de \textbf{3 columnas y 4 filas}. Columnas: \emph{Microsoft Word}, \emph{Google Docs}, \emph{LibreOffice Writer}. Filas: \emph{(a) ¿Es gratis o cuesta?}, \emph{(b) ¿Se usa con internet o sin internet?}, \emph{(c) ¿Dónde guarda los archivos?}, \emph{(d) ¿Qué pasa si lo abro en otro computador?}. Llena las 12 casillas con tu intuición. Si no sabes alguna, escribe \emph{``no sé''}. Después comparas con la sistematización.
\end{softbox}

\checkbox Dibujé la tabla 3x4 en el cuaderno.\\
\checkbox Llené las 12 casillas con mi intuición.\\
\checkbox No miré la siguiente actividad ni busqué en internet.

\begin{infoband}{milcVerde}
\textbf{Lo que va al cuaderno (Actividad 1 · IDENTIFICA).} Título: «Actividad 1 --- Comparando 3 procesadores». \textbf{Formato:} tabla 3 columnas × 4 filas. \textbf{Extensión:} 12 casillas llenas. \textbf{Sabes que terminaste cuando} tienes la tabla completa, aunque haya casillas con \emph{``no sé''}.
\end{infoband}

\puente{Después aprendes los 5 pasos fundamentales: crear, escribir, guardar, nombrar, ubicar.}

\begin{titlebox}{milcTurquesa}
Fase 2 · Sistematización con pensamiento computacional
\end{titlebox}

\begin{softbox}{milcTurquesa}{milcGris}{Cuatro pilares aplicados al tema}
Ahora compara tu tabla con la realidad. Vas a descubrir que cada procesador tiene ventajas distintas, y eso te ayuda a decidir cuál usar según la situación.
\end{softbox}

\small
\begin{tabularx}{\textwidth}{>{\bfseries\RaggedRight\arraybackslash}p{3.6cm}Y}
\rowcolor{milcTurquesa!90}\color{white}Pilar & \color{white}\bfseries Aplicación al tema\\
1. Descomposición & \textbf{Microsoft Word.} \emph{Costo:} cuesta (parte de Office 365), pero los colegios tienen licencia gratis para estudiantes. \emph{Internet:} se usa SIN internet. \emph{Guardado:} en tu disco duro (o en OneDrive si lo conectas). \emph{Otros equipos:} si llevas tu archivo en USB y abres en otro Word, todo se ve igual. \emph{Cuándo conviene:} para trabajos formales, tareas serias del colegio, cuando necesitas formato sofisticado.\\
2. Reconocimiento de patrones & \textbf{Google Docs.} \emph{Costo:} gratis (solo necesitas cuenta de Google). \emph{Internet:} se usa CON internet (también funciona algo sin internet, pero limitado). \emph{Guardado:} en la nube (Google Drive); puedes verlo desde cualquier dispositivo con tu cuenta. \emph{Otros equipos:} entras a docs.google.com en cualquier computador, te abre tu archivo. \emph{Cuándo conviene:} para trabajos en grupo (varios escribiendo a la vez), para no perder archivos al cambiar de computador, para colegio.\\
3. Abstracción & \textbf{LibreOffice Writer.} \emph{Costo:} gratis (descarga la suite). \emph{Internet:} se usa SIN internet. \emph{Guardado:} en tu disco. \emph{Otros equipos:} si la persona también tiene LibreOffice, se ve igual. Si la abren en Word, casi siempre se ve bien pero puede haber detalles distintos. \emph{Cuándo conviene:} si no tienes Word y no quieres pagar; en computadores Linux; en casa donde nadie tiene Office.\\
4. Algoritmo conceptual & \textbf{Los 5 pasos universales para crear y guardar un documento.} \textbf{(1) Crear:} abrir el programa (Word, Docs, etc.). \textbf{(2) Escribir:} el contenido en orden (título, párrafos). \textbf{(3) Guardar:} pulsar \texttt{Ctrl+G} (Guardar) o \texttt{Ctrl+S} en inglés (de \emph{Save}). \textbf{(4) Nombrar bien:} aplicar las 5 reglas que aprendiste en S7 (sin espacios, sin acentos, claro, con fecha si aplica). \textbf{(5) Ubicar en carpeta:} guardar en la carpeta correcta de tu árbol (no en el escritorio amontonado).\\
\end{tabularx}
\normalsize

\begin{drawbox}{milcTurquesa}{3 procesadores + 5 pasos universales}
\textbf{Procesadores:} Word (cuesta, sin internet) · Docs (gratis, con internet, en nube) · LibreOffice (gratis, sin internet). \\[1mm]
\textbf{5 pasos:} crear · escribir · guardar (Ctrl+G) · nombrar bien · ubicar en carpeta.
\end{drawbox}

\begin{softbox}{milcMostaza}{milcGris}{Errores comunes y su corrección}
\textbf{Errores que se cometen al empezar con procesadores.} (1) \emph{``No guardé nada hasta el final''} --- y el computador se apaga: pierdes todo. La regla es guardar cada 5 minutos. (2) \emph{``Lo guardé en el escritorio''} --- y a la semana no lo encuentras. La regla es ubicarlo en la carpeta correcta. (3) \emph{``Le puse de nombre tarea1.docx''} --- y después tienes tarea2, tarea3, tarea4. La regla es nombre claro con fecha. (4) \emph{``Hice todo el formato en Word y lo abrí en Docs y se desconfiguró''} --- normal. Para evitar: termina en el procesador en el que vas a entregar. (5) \emph{``Cerré sin guardar y dije sí a no guardar''} --- recuperación es difícil. Lee siempre los avisos antes de cerrar.
\end{softbox}

\begin{infoband}{milcTurquesa}
\textbf{Lo que va al cuaderno (Actividad 2 · EXPLICA).} Título: «Actividad 2 --- Los 5 pasos para escribir digital». \textbf{Formato:} lista numerada del 1 al 5. Cada paso: nombre + cómo se hace + atajo de teclado si aplica. \textbf{Extensión:} 5 líneas. \textbf{Sabes que terminaste cuando} podrías hacerle un tutorial verbal a tu mamá de cómo guardar un documento.
\end{infoband}

\puente{Con los pasos claros, eliges tu tema y produces tu documento digital de 3 párrafos.}

\begin{titlebox}{milcMagenta}
Fase 3 · Praxis con pensamiento computacional
\end{titlebox}

\begin{softbox}{milcMagenta}{milcGris}{Construyendo el producto con los cuatro pilares}
\textbf{Actividad 3 · APLICA --- Tu primer documento digital de 3 párrafos} (30 min · individual). Vas a crear tu \textbf{primer documento digital}. Si tienes computador en clase, lo haces ahí. Si no, lo planificas en el cuaderno y lo creas en casa esta semana. Tema libre: pueblo, tradición familiar, mascota, equipo de fútbol, materia favorita, lo que tú quieras.
\end{softbox}

\small
\begin{tabularx}{\textwidth}{>{\bfseries\RaggedRight\arraybackslash}p{3.6cm}Y}
\rowcolor{milcMagenta!90}\color{white}Pilar & \color{white}\bfseries Aplicación al producto\\
Descomposición & \textbf{Paso 1 --- Decide el tema y el procesador.} En el cuaderno escribe: \emph{``Tema: ...''} y \emph{``Procesador que voy a usar: Word / Docs / LibreOffice''}. Si en el colegio hay Word, úsalo. Si tienes cuenta Google, Docs también funciona.\\
Patrones & \textbf{Paso 2 --- Planifica los 3 párrafos en el cuaderno PRIMERO.} \textbf{Doña Mercedes lo dice}: \emph{antes de escribir, pensar}. En 3 líneas anota la idea principal de cada párrafo. Ejemplo: \emph{``Párrafo 1 (introducción): voy a hablar de mi gato Murcio. Párrafo 2 (desarrollo): cuento sus 3 características raras y una anécdota. Párrafo 3 (cierre): qué he aprendido de él.''}.\\
Abstracción & \textbf{Paso 3 --- Crea el documento.} Abre el procesador. Aparece una página en blanco. Escribe el título arriba (centrado): \emph{``Mi primer texto digital --- [tu tema]''}. Salta 2 líneas (presiona Enter 2 veces).\\
Algoritmo de armado & \textbf{Paso 4 --- Escribe los 3 párrafos.} Cada párrafo debe tener entre 3 y 5 líneas. \textbf{No copies de internet}: doña Mercedes te diría que la voz que vale es la tuya, no la de un sitio web. Si necesitas datos, los pones después con tus palabras. Escribe sin parar mucho; después revisas.\\
\end{tabularx}
\normalsize

\begin{drawbox}{milcMagenta}{Antes de mostrar el documento al profe}
\checkbox El tema lo escogí yo, no me lo impuso nadie. \\[1mm]
\checkbox Planifiqué los 3 párrafos en el cuaderno antes de escribir en digital. \\[1mm]
\checkbox El documento tiene título centrado arriba. \\[1mm]
\checkbox Los 3 párrafos tienen entre 3 y 5 líneas cada uno. \\[1mm]
\checkbox El documento está guardado con nombre que sigue las 5 reglas. \\[1mm]
\checkbox Está en la carpeta correcta de mi árbol (no en el escritorio).
\end{drawbox}

\begin{softbox}{milcMagenta}{milcGris}{Plantilla de trabajo}
\textbf{Plantilla del documento.} \textit{Arriba (centrado):} título de 1 línea. \textit{2 saltos.} \textit{Párrafo 1 (introducción):} 3-5 líneas, sobre qué vas a hablar. \textit{Párrafo 2 (desarrollo):} 3-5 líneas, lo principal del tema. \textit{Párrafo 3 (cierre):} 3-5 líneas, conclusión o aprendizaje. \textit{Guardado:} 2026-tema.docx en la carpeta de Tecnologia/Periodo-3.
\end{softbox}

\begin{infoband}{milcMagenta}
\textbf{Lo que va al cuaderno (Actividad 3 · APLICA).} Título: «Actividad 3 --- Mi primer documento digital». \textbf{Formato:} planificación de los 3 párrafos (3 líneas) + registro de creación (tema, procesador, ruta, tamaño) + si es posible, una captura impresa pegada o boceto del documento. \textbf{Extensión:} 1 página del cuaderno. \textbf{Sabes que terminaste cuando} tu documento abre desde su carpeta y se ve como lo planeaste.
\end{infoband}

\puente{El producto es ese documento guardado + el resumen en cuaderno.}

\begin{titlebox}{milcMostaza}
Producto de la sesión
\end{titlebox}
\textbf{Mi primer documento digital de 3 párrafos · guardado correctamente}.
\begin{softbox}{milcMostaza}{milcGris}{Criterios de calidad}
(1) El tema lo escogió el estudiante. (2) Hay 3 párrafos con sentido (introducción, desarrollo, cierre). (3) El nombre del archivo sigue las 5 reglas (sin espacios, sin acentos, claro, fecha). (4) Está en la carpeta correcta del árbol de organización. (5) El estudiante puede mostrar el archivo abriéndolo en clase o en casa.
\end{softbox}

\puente{Antes de salir, verificas que tu documento abre cuando lo das clic, que tiene 3 párrafos y que está en la carpeta correcta.}

\begin{titlebox}{milcVino}
Fase 4 · Evaluación liberadora — cinco dimensiones
\end{titlebox}

\begin{softbox}{milcVino}{milcGris}{Sentido de la evaluación}
Evaluar no es solo revisar una nota. Es reconocer qué aprendí, qué cambió en mí, qué emociones manejé, cómo me relaciono con la ciudad y la comunidad, y qué le dejaré a quien viene detrás.
\end{softbox}

\textbf{Mi reflexión liberadora en cinco dimensiones.} Completa cada frase iniciada:

\small
\begin{tabularx}{\textwidth}{>{\bfseries\RaggedRight\arraybackslash}p{3.4cm}Y>{\RaggedRight\arraybackslash}p{4.6cm}}
\rowcolor{milcVino}\color{white}Dimensión & \color{white}\bfseries Mi respuesta (frame starter) & \color{white}\bfseries Algo que descubrí\\
1. Personal & Lo que más cambió en mí fue \linead{6.5cm} & \linead{4.2cm}\\
2. Emocional & La emoción más fuerte fue \linead{2.5cm} y la regulé \linead{3cm} & \linead{4.2cm}\\
3. Ciudadana & En democracia esto importa porque \linead{6cm} & \linead{4.2cm}\\
4. Local (Cartago) & En mi barrio sería útil para \linead{6.5cm} & \linead{4.2cm}\\
5. Intergeneracional & A mi abuela/abuelo le diría \linead{6.5cm} & \linead{4.2cm}\\
\end{tabularx}
\normalsize

\begin{infoband}{milcVino}
\textbf{Para la próxima sustentación.} Vuelve a esta tabla en seis meses. Verás que las respuestas cambian — no porque lo aprendido cambie, sino porque tú cambias. La evaluación liberadora es una conversación contigo a través del tiempo.
\end{infoband}


\puente{Cerramos con 3 voces que te ayudan a entender por qué escribir bien sigue siendo un acto de respeto al lector.}

\begin{titlebox}{milcGradeDark}
Triángulo de pensamiento · cierre filosófico
\end{titlebox}

\begin{softbox}{milcVino}{milcGris}{Dussel · lente del nosotros}
\textit{````Escribir con tu voz, no con voz prestada. Esa es la primera disciplina del que aprende.''''}

Cuando empiezas a escribir digital, la tentación más grande es \textbf{copiar de internet}: hay millones de textos disponibles. Pero \emph{copiar no te enseña a escribir}. Doña Mercedes lo sabía: el alumno que copia se siente seguro hoy y queda sin voz mañana. El alumno que escribe lo suyo (aunque tenga errores) está construyendo su propia voz. La voz propia es lo único que de verdad importa al final.

\textbf{¿Estoy escribiendo con mi voz o con voces prestadas que aún no sé reconocer como mías?}
\end{softbox}
\linead{\linewidth}\\[1mm]
\linead{\linewidth}

\begin{softbox}{milcMostaza}{milcGris}{Estoicismo · Marco Aurelio (emperador que escribía sus reflexiones diarias) · cuidado interior}
\textit{````Escribir es pensar despacio. Hablar es pensar rápido. Los dos son necesarios; el lento construye más.''''}

Hablamos todo el día sin pensar mucho. \textbf{Escribir obliga a frenar}: a elegir palabra exacta, a ordenar la idea, a corregir. Por eso escribir te hace mejor pensador, aunque al principio sea difícil. \emph{Marco Aurelio escribió sus reflexiones más profundas en cartas y diarios}. Tú con un documento de 3 párrafos estás haciendo lo mismo: pensar despacio. Ese hábito te acompañará.

\textbf{¿Estoy aprovechando la escritura para pensar despacio, o solo escribo por obligación?}
\end{softbox}
\linead{\linewidth}\\[1mm]
\linead{\linewidth}

\begin{softbox}{milcTurquesa}{milcGris}{Floridi · lente de la infoesfera}
\textit{````El procesador de texto no cambia qué es escribir. Cambia el ritmo, la facilidad de revisar, la capacidad de compartir. Lo demás es lo mismo de siempre.''''}

Algunas personas creen que \textbf{escribir digital es ``moderno'' y escribir en cuaderno es ``viejo''}. Falso. \emph{Son herramientas distintas para el mismo oficio}. El cuaderno es íntimo y lento. El procesador es público y rápido. Los dos sirven, dependiendo de qué quieres escribir. La maestra rural del campo y la programadora de Silicon Valley hacen lo mismo: escribir para ser entendidas.

\textbf{¿Qué cosas escribo mejor en cuaderno y qué cosas escribo mejor en digital?}
\end{softbox}
\linead{\linewidth}\\[1mm]
\linead{\linewidth}

\begin{titlebox}{milcVino}
Compromiso final
\end{titlebox}

\small
\begin{tabularx}{\textwidth}{>{\RaggedRight\arraybackslash}p{3.0cm}YYY}
\rowcolor{milcVino}\color{white}\bfseries Criterio & \color{white}\bfseries Lo logré & \color{white}\bfseries En proceso & \color{white}\bfseries Necesito apoyo\\
Comprensión del tema & Explico ideas centrales con mis palabras. & Reconozco ideas pero falta precisión. & Necesito repasar conceptos base.\\
Pensamiento computacional & Apliqué los 4 pilares al diseñar mi producto. & Apliqué algunos pilares. & Necesito releer la fase 2.\\
Aplicación y producto & Mi producto cumple los criterios y se puede revisar. & Producto funcional con ajustes pendientes. & Aún en construcción.\\
Reflexión 5 dimensiones & Respondí cada dimensión con honestidad. & Respondí algunas con detalle. & Mis respuestas son muy generales.\\
Triángulo de pensamiento & Conecté las 3 voces con mi trabajo real. & Conecté al menos una. & No relaciono las voces con mi proceso.\\
\end{tabularx}
\normalsize

\textbf{Hoy entendí que:}\\[1mm]
\linead{\linewidth}\\[1mm]
\linead{\linewidth}

\textbf{Mi compromiso para la próxima sesión es:}\\[1mm]
\linead{\linewidth}\\[1mm]
\linead{\linewidth}

\begin{softbox}{milcMostaza}{milcGris}{Cita de despedida · Marco Aurelio}
\textit{``El obstáculo en el camino se vuelve el camino. Lo que estorba al camino se vuelve camino.''}\\[1mm]
Cada error de hoy, bien observado, se convierte en pieza del camino que pisarás mañana.
\end{softbox}

\small
\begin{tabularx}{\textwidth}{>{\bfseries}p{3.6cm}Y}
\rowcolor{milcGradeDark!12}Nombre completo: & \linead{10cm}\\
Fecha: & \linead{4cm} \quad Curso/grupo: \linead{4cm}\\
Mi firma: & \linead{9cm}\\
\end{tabularx}
\normalsize

\begin{infoband}{milcGradeDark}
\textbf{Para la próxima sesión.} Esta semana, abre tu primer documento al menos 2 veces más para releerlo y mejorarlo. Cada vez que lo releas, vas a notar algo: una palabra mejor, un párrafo más claro, un error que corregir. La próxima clase aprendes \textbf{formato básico}: tipos de letra, tamaños, alineación, color. Vas a transformar tu documento de \emph{texto plano} a \emph{texto con presencia}. Hoy escribiste; la próxima vez presentas.
\end{infoband}

\end{document}
